% !TeX program = lualatex
% Definition and sources only; this file is not an LLM solution attempt.
\documentclass[11pt,a4paper]{article}
\usepackage[margin=25mm]{geometry}
\usepackage{fontspec}
\setmainfont{lmroman10-regular.otf}[BoldFont=lmroman10-bold.otf,
 ItalicFont=lmroman10-italic.otf,BoldItalicFont=lmroman10-bolditalic.otf]
\usepackage{amsmath,amssymb,mathtools}
\usepackage{unicode-math}
\setmathfont{latinmodern-math.otf}
\AtBeginDocument{\let\setminus\smallsetminus}
\usepackage{xurl,hyperref}
\hypersetup{colorlinks=true,urlcolor=blue}
\setcounter{secnumdepth}{0}
\setlength{\parindent}{0pt}
\setlength{\parskip}{5pt}
\setlength{\emergencystretch}{3em}
\begin{document}
\section*{Research Problems in Function Theory — Problem 1.39}\label{statement}

\subsection{Definitions and mathematical statement}
For a nonconstant meromorphic function $f$, let $n(t,a)$ count solutions of $f(z)=a$ in $|z|\le t$, with multiplicity; $a=\infty$ means poles. Put
\[
N(r,a)=\int_0^r\frac{n(t,a)-n(0,a)}t\,dt+n(0,a)\log r,
\quad T(r,f)=N(r,\infty)+\frac1{2\pi}\int_0^{2\pi}\log^+|f(re^{i\theta})|\,d\theta,
\]
where $\log^+x=\max(0,\log x)$. At a boundary radius $R\in\{1,\infty\}$ with $T(r,f)\to\infty$, define the deficiency
$\delta(a,f)=1-\limsup_{r\to R}N(r,a)/T(r,f)$. A deficient value has positive deficiency.
Work in the unit disc $\mathbb D=\{z:|z|<1\}$ and define
\[
\alpha(f)=\limsup_{r\uparrow1}\frac{T(r,f)}{\log(1/(1-r))}.
\]
Consider only functions for which this number is finite. Determine the two sharp suprema of $\alpha(f)$ under the following respective hypotheses:
\begin{enumerate}
\item $f$ is holomorphic and nowhere zero in $\mathbb D$, and $f'(z)\ne1$ everywhere;
\item $f$ is meromorphic and has no zeros in $\mathbb D$, and $f'(z)\ne1$ wherever the derivative is finite.
\end{enumerate}
In the second class poles are allowed. The source gives universal upper bounds $2$ for the first class and $7$ for the second. It asks whether $2$ is optimal and what the optimal bound in the meromorphic class is. Constants can be included without changing a positive sharp supremum. The finite-growth assumption is retained rather than inferred from the omitted-value conditions.
\subsection{Short English statement}
Find the largest possible logarithmic growth for zero-free disc functions whose derivative omits one, distinguishing holomorphic from meromorphic functions.
\subsection{Sources}
\begin{itemize}
\item[S1] ULAM.AI and the UnsolvedMath Contributors, supplied problems.json, catalogue entry, AMR Open Problem Lists. Catalogue identity and source transcription; CC BY 4.0.
\url{https://huggingface.co/datasets/ulamai/UnsolvedMath}.
\item[S2] Hayman - Research Problems in Function Theory (2018), Problem 1.39. 
\url{https://arxiv.org/abs/1809.07200}.
\end{itemize}
\end{document}
